\chapter{Definición del problema}
\label{defProblema}  % para poder enlazar a este capítulo

\section{Universidad: contexto de la labor docente}

Ejemplo de un texto marcado como literal de una fuente citada: 

Las funciones básicas de la enseñanza universitaria, tal y como se concibe en la actualidad, fueron establecidas por José Ortega y
Gasset en su trabajo \textquote{Misión de la Universidad}~\cite{OrtegayGasset1982}.

Ejemplo de estilo para citar un bloque de texto: 

Las funciones de la Universidad propuestas por Ortega y Gasset se encuentran recogidas en la Ley Orgánica de Universidades (LOU) que establece:

\begin{displayquote}[Ley Orgánica de Universidades~\cite{LOU}]
Artículo 1. Funciones de la Universidad:
\begin{enumerate}
\item La Universidad realiza el servicio público de la educación superior, mediante la investigación, la
docencia y el estudio.
\item Son funciones de la Universidad al servicio de la sociedad:
    \begin{enumerate}
    \item La creación, desarrollo, transmisión y crítica de la ciencia, de la técnica y de la cultura.
    \item La preparación para el ejercicio de actividades profesionales que exijan la aplicación de
    conocimientos y métodos científicos y para la creación artística.
    \item La difusión, la valorización y la transferencia del conocimiento al servicio de la cultura, de
    la calidad de vida, y del desarrollo económico.
    \item La difusión del conocimiento y la cultura a través de la extensión universitaria y la formación a lo largo
    de toda la vida.
    \end{enumerate}
\end{enumerate}

\end{displayquote}

Ejemplo de figura mencionada en el texto: Según los datos del Ministerio de Educación, Cultura y Deporte publicados en su memoria del curso 2015-2016~\cite{DatosMECD}, que es la más actual disponible en la fecha de realización de este trabajo, los datos del número de alumnos apuntan a una estabilización en los últimos años. Sin embargo, si se observan los datos de forma relativa a la población siguen mostrando un crecimiento importante. En la figura~\ref{fig.0201EvolucionMatriculados}, se pueden observar estas tendencias.

\begin{figure}[hbtp]
	\centering
	\includegraphics[width=1.0\textwidth]{0201EvolucionMatriculados2016.png}
	\caption{Evolución de la población de 18 a 24 años, de los matriculados en Grado, 1 er y 2o ciclo y Máster y de la
tasa neta de escolarización en educación universitaria~\cite{DatosMECD}}
	\label{fig.0201EvolucionMatriculados}
\end{figure}

% Nota solo para el autor de la plantilla: XXX Poner solo una que sea PDF pero cite bien la fuente (para que se vea que es mejor vectorial, o una de cada ejemplo captura pantalla justificación mapa bits)